\section{Distribuição Normal}

A distribuição normal foi introduzida em 1733 por Abraham DeMoivre para aproximar
probabilidades de variáveis binomiais quando o parâmetro $n$ é grande.
Essa distribuição possui grande importância na Teoria de Probabilidades e na Estatística,
pois suas propriedades analíticas permitem o tratamento adequado de inúmeros fenômenos

\begin{def:normal}
A variável contínua $X$ tem distribuição normal com parâmetros $\mu$ e $\sigma^2$ se
sua densidade for
\[
    f(x) = \frac{1}{\sqrt{2\pi\sigma^2}} e^{\frac{-(x - \mu)^2}{2\sigma^2}}.
\]
Uma variável normal será padrão caso $\mu = 0$ e $\sigma^2 = 1$.
\end{def:normal}

\begin{lem:integral exponencial quadrada}
\label{lem:integral exponencial quadrada}
A integral definida de $f(x) = e^{\frac{-x^2}{2}}$ avaliada em $(-\infty, \infty)$ é
$\sqrt{2\pi}$.
\[
    \int_{-\infty}^\infty e^{\frac{-x^2}{2}} = \sqrt{2\pi}
\]
\end{lem:integral exponencial quadrada}
\begin{proof}
Equivalentemente, podemos provar que
\[
    \left(\int_{-\infty}^{\infty}e^{\frac{-u^2}{2}}\,du\right) \left(\int_{-\infty}^{\infty}e^{\frac{-t^2}{2}}\,dt\right) = 2\pi.
\]
Espressando o problema como uma integral dupla, tal que
\[
    \int_{-\infty}^{\infty} \int_{-\infty}^{\infty} e^{\frac{-(t+u)^2}{2}}\,dt\,du,
\]
e colocando em termos de coordenadas polares, isto é,
\[
    r^2 = u^2 + t^2, \quad t = r \cos \theta, \quad u = r \cos \theta,
\]
obtemos
\[
    \int_0^\infty \int_0^{2\pi} re^{\frac{-r^2}{2}}\,dr\,d \theta =
    2 \pi \int_0^\infty re^{\frac{-r^2}{2}}\,dr,
\]
essa que pode ser resolvida via substituição, resultando em
\[
    - 2 \pi \left. e^{\frac{-r^2}{2}} \right|_0^\infty = 2 \pi.
\]
\end{proof}

\begin{prop:densidade da normal}
A função
\[
    f(x) = \frac{1}{\sqrt{2\pi\sigma^2}} e^{\frac{-(x - \mu)^2}{2\sigma^2}}
\]
é uma função de densidade.
\end{prop:densidade da normal}
\begin{proof}
Para tanto, devemos mostrar que
\[
    \frac{1}{\sqrt{2\pi\sigma^2}} \int_{-\infty}^{\infty}e^{\frac{-(x - \mu)^2}{2\sigma^2}}\,dx = 1
\]
Realizemos a substituição
\[
    z = \frac{x - \mu}{\sigma},
\]
de modo que
\[
    \frac{1}{\sqrt{2\pi\sigma^2}} \int_{-\infty}^{\infty}e^{\frac{-(x - \mu)^2}{2\sigma^2}}\,dx =
    \frac{1}{\sqrt{2\pi}} \int_{-\infty}^{\infty}e^{\frac{-z^2}{2}}\,dz = 1.
\]
Resta concluir que
\[
    \int_{-\infty}^{\infty}e^{\frac{-z^2}{2}}\,dz = \sqrt{2\pi},
\]
o que é garantido pelo Lema \ref{lem:integral exponencial quadrada}.
\end{proof}

\begin{thm:tranformacao afim da normal}
Se $X \sim \mathrm{Normal}(\mu, \sigma^2)$ e $a$ e $b$ são constantes, então
a variável $aX + b$ é normalmente distribuída com parâmetros $a\mu + b$ e $a^2\sigma^2$.
\end{thm:tranformacao afim da normal}
\begin{proof}
Se a distribuição de $X$ é dada por
\[
    \int_{-\infty}^\infty \frac{1}{\sqrt{2\pi\sigma^2}} e^{\frac{-(x - \mu)^2}{2\sigma^2}}
\]
e temos que
\[
    y = ax + b \Longrightarrow x = \frac{y - b}{\mu}
\]
então, substituindo $x$ na função de distribuição, obtemos
\[
    \frac{1}{\sqrt{2\pi a^2\sigma^2}} \int_{-\infty}^{\infty} e^{\frac{-(y - a\mu - b)^2}{2a^2\sigma^2}}\,dy.
\]
Logo, a expressão anterior trata-se da distribuição de uma variável $Y = aX + b$ tal que
$Y \sim \mathrm{Normal}(a\mu + b, a^2\sigma^2)$.
\end{proof}

\begin{cor:normal e normal padrao}
\label{cor:normal e normal padrao}
Se $X \sim \mathrm{Normal}(\mu, \sigma^2)$ e $Z = \frac{X - \mu}{\sigma}$, então
$Z \sim \mathrm{Normal}(0, 1)$.
\end{cor:normal e normal padrao}
\begin{proof}
Se $\mu_z$ e $\sigma^2_z$ são os parâmetros da distribuição de $Z$, então
\[
    \mu_z = \frac{1}{\sigma} \cdot \mu - \frac{\mu}{\sigma} = 0 \quad \text{e} \quad
    \sigma_z^2 = \left(\frac{1}{\sigma}\right)^2 \cdot \sigma^2 = 1
\]
e
\end{proof}

\begin{thm:parametros da normal}
Se $X \sim \mathrm{Normal}(\mu, \sigma^2)$, então
\[
    \ev{X} = \mu \quad \text{ e } \quad \var{X} = \sigma^2
\]
\end{thm:parametros da normal}
\begin{proof}
Se $Z \sim \mathrm{Normal}(0, 1)$, então
\[
    \ev{Z} = \frac{1}{\sqrt{2\pi}} \int_{-\infty}^\infty z e^{\frac{-z^2}{2}}.
\]
Resolvendo por substituição $u = \frac{-z^2}{2}$, obtemos
\[
    \ev{Z} = \frac{1}{\sqrt{2\pi}} \left. e^{\frac{-z^2}{2}} \right|_{-\infty}^\infty = 0.
\]
Por outro lado,
\[
    \var{Z} = \frac{1}{\sqrt{2\pi}} \int_{-\infty}^\infty z^2 e^{\frac{-z^2}{2}},
\]
a qual pode ser resolvida usando-se integração por partes com $u = z$ e
$dv = -ze^{\frac{-z^2}{2}}dz$, obtendo $v = e^{\frac{-z^2}{2}}dz$ e
\[
    \var{Z} = -\frac{1}{\sqrt{2\pi}} \left( \left. z^2 e^{\frac{-z^2}{2}} \right|_{-\infty}^\infty -
    \int_{-\infty}^{\infty} e^{\frac{-z^2}{2}}\,dz \right).
\]
Utilizando o Lema \ref{lem:integral exponencial quadrada}, encontramos que
\[
    \var{Z} = \frac{1}{\sqrt{2\pi}} \cdot \sqrt{2\pi} = 1.
\]
Então conclui-se que o Teorema vale para uma variável normal padrão.
Se $X \sim \mathrm{Normal}(\mu, \sigma^2)$ e $Z = \frac{X - \mu}{\sigma}$, então, pelo Corolário
\ref{cor:normal e normal padrao}, $Z$ é normal padrão.
Pelos Teoremas \ref{thm:propriedades esperanca} e \ref{prop:propriedades variancia}, temos
que
\[
    \ev{X} = \ev{\sigma Z + \mu} = \sigma \ev{Z} + \mu \quad \text{e} \quad
    \var{X} = \var{\sigma Z + \mu} = \sigma^2 \var{Z}.
\]
Mas como $\ev{Z} = 0$ e $\var{Z} = 1$, então
\[
    \ev{X} = \mu \quad \text{e} \quad \var{X} = \sigma^2
\]
\end{proof}

\begin{prop:simetria densidade normal}
\label{prop:simetria densidade normal}
A função de densidade $f(x)$ de uma variável normal é uma função par, isto é,
\[
    f(-x) = f(x)
\]
\end{prop:simetria densidade normal}
\begin{proof}
	Seja $X \sim \mathrm{Normal}(\mu, \sigma^2)$.
	Uma vez que
	\[
	    | - x - \mu | = | x - \mu | \Longrightarrow (-x-\mu)^2 = (x - \mu)^2,
	\]
	então
	\[
	    f(-x) = \frac{1}{\sqrt{2\pi\sigma^2}} e^{\frac{-(-x-\mu)^2}{2\sigma^2}} =
		\frac{1}{\sqrt{2\pi\sigma^2}} e^{\frac{-(x-\mu)^2}{2\sigma^2}} = f(x).
	\]
\end{proof}

\begin{thm:simetria distribuicao normal}
Se $F(x)$ for a função de distribuição de uma variável normal, então
\[
    F(-x) = 1 - F(x)
\]
\end{thm:simetria distribuicao normal}
\begin{proof}
Com efeito,
\[
    F(-x) = \int_{-\infty}^{-x} f(u)\,du.
\]
onde $f$ é a função de densidade da normal.
Realizando a substituição $u = -t$, então $du = -dt$.
Além disso, temos que $u \to -\infty$ implica em $t \to \infty$, enquanto que
$u = -x$ nos dá $t = x$.
Obtemos, por consequência,
\[
    F(-x) = - \int_{\infty}^{x} f(-t)\,du.
\]
Pelas propriedades de integrais,
\[
    F(-x) = \int_{x}^{\infty} f(-t)\,du,
\]
e como $f$ é par pela Proposição \ref{prop:simetria densidade normal}, então
\[
    F(-x) = \int_{x}^{\infty} f(t)\,du.
\]
No entanto,
\begin{gather*}
	\int_{-\infty}^x f(t)\,dt + \int_{x}^\infty  = 1 \\
	\int_{x}^\infty f(t)\,dt = 1 - \int_{-\infty}^x f(t)\,dt\\
	\int_{x}^\infty f(t)\,dt = 1 - F(x)\\
\end{gather*}
mostrando, enfim, que
\[
    F(-x) = 1 - F(x)
\]
\end{proof}

\begin{thm:demoivre-laplace}[de DeMoivre-Laplace]
Se $X \sim \mathrm{Binomial}(n, p)$, então para quaisquer reais $a$ e $b$ com $a < b$ é
satisfetio que
\[
    \lim_{n \to \infty} P\left(a \le \frac{X - np}{\sqrt{np(1-p)}} \le b\right) = \Phi(b) - \Phi(a)
\]
\end{thm:demoivre-laplace}
